\documentclass[10pt]{article}

\usepackage[utf8]{inputenc}

\usepackage[T1]{fontenc}

\usepackage{amsmath}

\usepackage{amsfonts}

\usepackage{amssymb}

\usepackage[version=4]{mhchem}

\usepackage{stmaryrd}

\usepackage{bbold}

\usepackage{graphicx}

\usepackage[export]{adjustbox}

\graphicspath{ {./images/} }



\begin{document}

оно получено Ж. Лагранжем (Ј. Lagrange, 1760) и истолковано ЖК. Мёнье (J. Meusnier) как условие равенства нулю средней кривизны поверхности $\boldsymbol{z}=\boldsymbol{z}(x, y)$, частные интегралы найдены Г. Монжем (G. Monge). Систематические исследования Ә.-Л. у. проведены С. М. Бернштейном, который показал, что Э.-Л. у. является квазилинейным эллиптич. уравнением рода $p=2$, вследствие чего решения Ә.-Л. у. обладают рядом свойств, резко отличающих их от решений линейных уравнений. К таким свойствам, напр., относятся устранимость изолированных особых точек решения без априорного предположения об ограниченности решения в окрестности особой точки, принцип максимума, имеющий место при тех же условиях, невозможность равномерной априорной оценки $z(x, y)$ в любой компактной подобласти круга через значения $z$ в центре круга (т, е. отсутствие точного аналога неравенства Гарнака), факты, относящиеся к Дирихле задаче, отсутствие нелинейного решения Э.-Л.У., определенното над всей плоскостью (Бернштейна теорема) и т. Д.



Э.-Л. у. обобщается по размерности: для минимальной типерповерхности $z=z\left(x_{1}, \ldots, x_{n}\right)$ в $\mathbb{R}^{n+1}$ соответствующее уравнение имеет вид



$\sum_{i=1}^{n} \frac{\partial}{\partial x_{i}}\left(\frac{\partial z}{\partial x_{i}} / \sqrt{1+|\nabla z|^{2}}\right)=0, \nabla z=\left(\frac{\partial z}{\partial x_{1}}, \ldots, \frac{\partial z}{\partial x_{n}}\right)$.



Для этого уравнения $(n \geqslant 3)$ исследована разрешимость задачи Дирихле, доказана устранимость особенностей решения, если они сосредоточены внутри области на множестве нулевой $(n-1)$-мерной меры Хаусдорфа, показана справедливость теоремы Бернштейна для $n \leqslant 7$ и построены контрпримеры для $n \geqslant 8$.



ЭЙЛЕРА - МАКЛОРЕНА ФОРМУЛА - формула суммирования, связывающая частные суммы ряда с интегралом и производными его общего члена:



\[

\begin{gathered}

\sum_{k=p}^{m-11} \varphi(k)=\int_{p}^{m} \varphi(t) d t+ \\

+\sum_{v=1}^{n-1} \frac{B_{v}}{v !}\left\{\varphi^{(v-1)}(m)-\varphi^{(v-1)}(p)\right\}+R_{n},

\end{gathered}

\]



где $B_{v}$ - Бернулли числа, $R_{n}$ - остаточный член. С помощью Бернулли многочленов $B_{n}(t), B_{n}(0)=B_{n}$ остаточный член записывается в виде:



\[

R_{n}=-\frac{1}{n !} \int_{0}^{1}\left[B_{n}(t)--B_{n}\right] \sum_{k=p}^{m-1} \varphi^{(n)}(k+1-t) d t

\]



Для $n=2 s$ остаточный член $R_{2 s}$ может быть представлен с использованием чисел Бернулли:



\[

R_{2 s}=\frac{B_{2 s}}{(2 s) !} \sum_{k=p}^{m-1} \varphi^{(2 s)}(k+\theta), \quad 0<\theta<1 .

\]



Если производные $\varphi^{(2 s)}(t)$ и $\varphi^{(2 s+1)}(t)$ имеют одинаковые знаки и не меняют знака на $[p, m]$, то



\[

R_{2 s}=\theta \frac{B_{2 s}}{(2 s) !}\left[\varphi^{(2 s-1)}(m)-\varphi^{(2 s-1)}(p)\right], 0 \leqslant \theta \leqslant 1 .

\]



Если, кроме того,



\[

\lim _{x \rightarrow \infty} \varphi^{(2 s-1)}(x)=0 \text {, }

\]



то Э.-М. Ф. может быть записана в виде:



\[

\begin{gathered}

\sum_{k=p}^{m-1} \varphi(k)=c+\int_{p}^{m} \varphi(t) d t+ \\

+\sum_{k=1}^{2 \varsigma-2} \frac{B_{k}}{k !} \varphi^{(k-1)}(m)+\theta \frac{B_{2 s}}{(2 s) !} \varphi^{(2 s-1)}(m), 0<\theta<1 .

\end{gathered}

\]



В такой форме Э.-М. ф. применяется, напр., при выводе Стирлинга бормулы. В этом случае $\varphi(x)=\ln x$ и c - Әйлера постоянная. Имеются обобщения Э. - М. ф. на случай кратных сумм.



Э. -М. ф. применяется для приближенного вычисления определенных интегралов, для шсследования сходи- мості рядов, для вычисления сумм п для разложения функций в ряд Тейлора. Напр., при $m=1, p=0, n=$ $=2 m+1, \varphi(x)=\cos \frac{x t}{2}$ Э.-М. ф. дает следующее выражение:



\[

\begin{gathered}

\frac{t}{2} \operatorname{ctg} \frac{t}{2}=\sum_{v=0}^{m}(-1)^{v} \frac{t^{2 v}}{(2 v) !} B_{2 v}+ \\

+\frac{(-1)^{m+1} t^{2 m+2}}{2 \sin \frac{t}{2}} \int_{0}^{1} \frac{B_{2 m+1}(t)}{(2 m+1) !} \sin \left(x-\frac{1}{2}\right) t d x .

\end{gathered}

\]



Э.-М. Ф. играет важную роль при изучении асимптотич. разложений, в теоретико-числовых оценках, в конечных разностей исчислении.



Э.-М. ф. иногда применяется в виде:



\[

\begin{gathered}

\sum_{0}^{n} \varphi_{n}(x)=\int_{0}^{n} \varphi(x) d x+\frac{1}{2}\left(\varphi_{0}+\varphi_{n}\right)+ \\

+\int_{0}^{n}\left(x-[x]-\frac{1}{2}\right) \varphi^{\prime}(x) d x .

\end{gathered}

\]



Э.-М. Ф. была впервые приведена Л. Эйлером [1] в виде:



\[

S=\int t d n+\alpha t+\beta \frac{d t}{d n}+\gamma \frac{d^{2} t}{d n^{2}}+\delta \frac{d^{3} t}{d n^{8}}+\varepsilon \frac{d^{4} t}{d n^{4}}+\ldots,

\]



где $S$ - сумма первых членов ряда с общим членом $t(n), S=t=0$ при $n=0$, а коэффициенты определяются рекуррентными соотношениями:



\[

\begin{aligned}

& \alpha=\frac{1}{2}, \beta=\frac{\lceil\alpha}{2 !}-\frac{1}{3 !}=\frac{1}{12}, \gamma=\frac{\beta}{2 !}-\frac{\alpha}{3 !}+\frac{1}{4 !}=0, \\

& \delta=\frac{\gamma}{2 !}-\frac{\beta}{3 !}+\frac{\alpha}{4 !}-\frac{1}{5 !}=-\frac{1}{720}, \varepsilon=0, \gamma=0, \ldots .

\end{aligned}

\]



Независимо формула была открыта позднее F. Маклореном [2].



Jum.: [1] E u l e r J., "Comment. Acad. sci. Imp. Petrop.»,



\includegraphics[max width=\textwidth, center]{2023_07_09_20c51d97fb68ef3e0f92g-1}

ряды, пер. с англ., М., 1951; [4] N ö r 1 u n d N̈. E., Vorlesungen über Differenzenrechnung, В., 1924; [5] $\Gamma$ е л ь ф о н д А. O., Псчисление конечных разностей, 3 изд., М., 1967. әффициентов $\Phi$ yрье ряда.



ЭЙЕРОВ КЈАСС - первое препятствие к построению сечения расслоения со слоем сфера, ассоции рованного с векторным расслоением. См. Характеристический класс.



ЭЙЛЕРОВА ХАРАКТЕРИСТИКА конечного клеточного комплекса $K$ - целое число



\[

\chi(K)=\sum_{k=0}^{\infty}(-1)^{k} \alpha_{k},

\]



гле $\alpha_{k}$ - число $k$-мерных клеток комплекса. Названа в честь Л. Эйлера (L. Euler), к-рый доказал в 1758, тто число вершин $\mathrm{B}$, ребер $\mathrm{P}$ и граней Г. выпуклого многогранника связаны формулой $\mathrm{B}-\mathrm{P}+\Gamma=2$. В неявном виде эта формула была известна еще Р. Декарту (R. Descartes, 1620). Оказывается, что



\[

\chi(K)=\sum_{k=0}^{\infty}(-1)^{k} p^{k}

\]



где $p^{k}$ есть $k$-мервое м у л а Э ӥ л е р а - П у а н к а р е). Э. х. комплекса $K$ является его гомологическим, гомотопитеским и топологическим инвариантом. В частности, Э. х. не зависит от способа разбиения пространства на клетки. Поэтому можно говорить, напр., об Э.х. произвольного компактного полиэдра, понимая под этим Э.х. какойнибудь его триангуляции. С другой стороны, формула Эїлера-Пуанкаре позволяет распространить понятие Э. х. на более широкий класс пространств и пар пространств, для к-рых правая часть формулы имеет смысл. Для произвольного поля $F$ справедлива обобценная формула Эӥлера - Пуанкаре, выражаюцая Э. х. через





\end{document}